\chapter{Las plantas son productoras}\label{ch:g6:plants-are-producers}

Pesa una maceta regada, \omterm{def:g1:plants-around-us:plant}{planta} en ella una bellota y vuelve al cabo de
unos años a pesar el roble joven: kilos de madera, corteza, \omterm{def:g1:plants-around-us:parts}{hoja} y \omterm{def:g1:plants-around-us:parts}{raíz}
--- mientras que la tierra de la maceta apenas ha perdido un gramo (el
\cref{ex:g4:what-plants-need:tree} soltó esta sorpresa por primera
vez). Kilos de árbol, ¿hechos de qué? Este capítulo deja clara la
respuesta --- y le pone nombre al oficio exclusivo de las \omterm{def:g1:plants-around-us:plant}{plantas}.

\section{Materia orgánica y materia mineral}

\begin{definition}[Materia mineral y materia orgánica]\label{def:g6:plants-are-producers:matter}
Para un biólogo, la materia es de dos clases:
\begin{itemize}
  \item la \emph{materia mineral}\index{materia mineral}: la materia
        del mundo no vivo --- el agua, los gases del aire, la roca, las
        \omterm{prop:g4:what-plants-need:needs}{sales minerales} disueltas de la
        \cref{prop:g4:what-plants-need:needs};
  \item la \emph{materia orgánica}\index{materia orgánica}: la materia
        fabricada por los \omterm{def:g6:exploring-our-environment:components}{seres vivos} y que forma sus cuerpos ---
        madera, \omterm{def:g1:plants-around-us:parts}{hoja}, \omterm{def:g3:bones-and-muscles:muscle}{músculo}, grasa, azúcar, lana. Se quema, se pudre
        y siempre se le puede rastrear un fabricante vivo.
\end{itemize}
\end{definition}

\begin{example}[Clasificar lo que hay en un bolsillo]\label{ex:g6:plants-are-producers:sorting}
Del bolsillo de un abrigo después de un paseo: un guijarro ---
mineral; una avellana --- orgánica (la fabricó un avellano); un botón
de plástico --- orgánico por su origen, curiosamente (los plásticos se
hacen a partir del petróleo, los restos transformados de \omterm{def:g6:exploring-our-environment:components}{seres vivos}
antiquísimos); una gota de lluvia --- mineral; un guante de lana ---
orgánico (una oveja fabricó la lana).
\end{example}

\section{El oficio de las productoras}

\begin{proposition}[Las plantas producen materia orgánica a partir de materia mineral]\label{prop:g6:plants-are-producers:produce}
Una \omterm{def:g1:plants-around-us:plant}{planta} verde, a la que solo se le dan ingredientes
\emph{minerales} --- agua y sales disueltas de la tierra, un gas tomado
del \emph{aire} --- y la \omterm{prop:g3:balanced-diet:jobs}{energía} de la \emph{\omterm{def:g6:exploring-our-environment:conditions}{luz}}, fabrica materia
\emph{orgánica}: su propio azúcar y, a partir de él, madera, \omterm{def:g1:plants-around-us:parts}{hoja} y
\omterm{def:g1:plants-around-us:parts}{raíz}. Esa fabricación es el oficio que la
\cref{prop:g4:what-plants-need:make} prometió explicar: las \omterm{def:g1:plants-around-us:plant}{plantas} son
\emph{productoras}\index{productor} --- las únicas fabricantes de
\omterm{def:g6:plants-are-producers:matter}{materia orgánica} nueva a partir de suministros puramente minerales. La
sustancia verde de las \omterm{def:g1:plants-around-us:parts}{hojas} es la maquinaria que atrapa la \omterm{def:g6:exploring-our-environment:conditions}{luz}; la
fabricación ocurre en la \omterm{def:g1:plants-around-us:parts}{hoja}, de día.
\end{proposition}
\admitted

\begin{omfigure}
\begin{tikzpicture}[scale=0.92]
  % leaf in the middle
  \draw[thick, omMeth, fill=omMeth!20] (5.2,1.5)
    .. controls (5.6,2.4) and (6.8,2.5) .. (7.2,1.6)
    .. controls (6.8,0.8) and (5.6,0.75) .. (5.2,1.5);
  \draw[omMeth!60!black] (5.35,1.5) -- (7.05,1.58);
  \node[font=\small\bfseries, omMeth!60!black] at (6.2,2.75)
    {la hoja verde, de día};
  % inputs
  \node[font=\small, align=center] (water) at (1.4,0.35)
    {agua y sales minerales\\ (de la tierra, por las raíces)};
  \node[font=\small, align=center] (gas) at (1.4,1.7)
    {un gas del aire};
  \node[font=\small, align=center] (light) at (1.4,3.1)
    {energía de la luz};
  \draw[->, thick, omDef] (water.east) -- (5.3,1.2);
  \draw[->, thick, omDef] (gas.east) -- (5.15,1.5);
  \draw[->, thick, omProp] (light.east) -- (5.5,2.05);
  % outputs
  \node[font=\small, align=center] (sugar) at (10.6,1.9)
    {\textbf{materia orgánica}:\\ azúcar y después madera,\\ hoja, raíz, semilla};
  \draw[->, very thick, omMeth!60!black] (7.3,1.7) -- (sugar.west);
  % oxygen small note
  \node[font=\footnotesize, align=center] (oxy) at (10.6,0.5)
    {(y oxígeno, soltado\\ al aire --- véase la observación)};
  \draw[->, omExo] (7.25,1.2) -- (oxy.west);
\end{tikzpicture}

{\small El oficio de las \omterm{def:g5:food-webs:roles}{productoras}: entran \omterm{def:g6:plants-are-producers:matter}{materia mineral} y \omterm{def:g6:exploring-our-environment:conditions}{luz},
sale \omterm{def:g6:plants-are-producers:matter}{materia orgánica}. Solo las \omterm{def:g1:plants-around-us:plant}{plantas} verdes (y sus parientes verdes)
llevan esta fabricación.}
\end{omfigure}

\begin{example}[Las pruebas, reunidas]\label{ex:g6:plants-are-producers:evidence}
Tres observaciones anteriores encajan ahora entre sí. La maceta pesada:
los kilos del roble no salieron de la tierra --- salieron sobre todo
del aire y del agua. El prisionero pálido del
\cref{ex:g4:what-plants-need:pale}: sin \omterm{def:g6:exploring-our-environment:conditions}{luz} no hay fabricación --- la
\omterm{def:g1:plants-around-us:plant}{planta} se muere de hambre a oscuras con la despensa de minerales llena.
La menta hambrienta de minerales del \cref{exo:g4:what-plants-need:11}:
las sales hacen falta --- pero como suministros pequeños, no como el
grueso. El grueso, del aire y del agua; la \omterm{prop:g3:balanced-diet:jobs}{energía}, de la \omterm{def:g6:exploring-our-environment:conditions}{luz}; el
condimento, de la tierra.
\end{example}

\begin{remark}[El intercambio de gases, en los dos sentidos]\label{rem:g6:plants-are-producers:gas}
El gas del aire que toma la \omterm{def:g1:plants-around-us:parts}{hoja} es el \emph{dióxido de carbono} ---
justo el gas que espira todo el que respira
(\cref{prop:g4:breathing:exchange}); y la fabricación suelta al aire
\emph{oxígeno} --- justo el gas que necesita todo el que respira. Las
\omterm{def:g1:plants-around-us:plant}{plantas} también respiran, de día y de noche, como todos los \omterm{def:g6:exploring-our-environment:components}{seres
vivos}; pero de día la fabricación de una \omterm{def:g1:plants-around-us:parts}{hoja} en funcionamiento supera
con mucho a su \omterm{def:g4:breathing:movements}{respiración}. El aire respirable del mundo es un
subproducto de las \omterm{def:g5:food-webs:roles}{productoras} --- una deuda más que todos los \omterm{def:g1:animals-around-us:animal}{animales}
tienen con lo verde.
\end{remark}

\section{Las productoras alimentan a todo el mundo --- y guardan para sí mismas}

\begin{proposition}[El primer eslabón, explicado]\label{prop:g6:plants-are-producers:firstlink}
La \cref{prop:g3:food-chains:start} decretó que toda \omterm{def:g3:food-chains:chain}{cadena
alimentaria} empieza por una \omterm{def:g1:plants-around-us:plant}{planta}; ahora el decreto tiene su razón.
Solo las \omterm{def:g5:food-webs:roles}{productoras} fabrican \omterm{def:g6:plants-are-producers:matter}{materia orgánica} a partir de \omterm{def:g6:plants-are-producers:matter}{materia
mineral}; todos los \omterm{def:g5:food-webs:roles}{consumidores} (\cref{def:g5:food-webs:roles}) solo
pueden \emph{transformar} la \omterm{def:g6:plants-are-producers:matter}{materia orgánica} que comen. Así que toda
la \omterm{def:g6:plants-are-producers:matter}{materia orgánica} de todos los cuerpos vivos --- conejo, zorro,
ratonero, tú --- se fabricó primero en una \omterm{def:g1:plants-around-us:parts}{hoja}, y la pirámide
alimentaria de la \cref{prop:g5:food-webs:pyramid} se apoya en el piso
verde porque ahí es por donde la materia entra en el mundo vivo.
\end{proposition}
\admitted

\begin{example}[La despensa propia de la planta]\label{ex:g6:plants-are-producers:pantry}
Lo que una \omterm{def:g1:plants-around-us:plant}{planta} fabrica de más para las necesidades del día, lo
guarda como reservas orgánicas: el \omterm{ex:g4:plants-without-seeds:bulbtuber}{tubérculo} de la patata y el \omterm{ex:g4:plants-without-seeds:bulbtuber}{bulbo}
del tulipán (\cref{ex:g4:plants-without-seeds:bulbtuber}), las mitades
gordas de la judía (\cref{def:g2:from-seed-to-plant:seed}), la \omterm{def:g1:plants-around-us:parts}{raíz}
gorda de la zanahoria (\cref{exo:g1:plants-around-us:10}), el dulzor de
la remolacha. Cada una es \omterm{def:g6:exploring-our-environment:conditions}{luz} de verano ingresada en forma de materia
--- y cada una es exactamente lo que cosechan los agricultores y lo que
se comen los comensales.
\end{example}

\begin{method}[Rastrear cualquier alimento hasta su hoja]\label{met:g6:plants-are-producers:trace}
Coge cualquier alimento de la mesa:
\begin{enumerate}
  \item si es materia vegetal --- el trigo del pan, la patata, la
        manzana --- es fabricación directa de una \omterm{def:g5:food-webs:roles}{productora}: nombra la
        \omterm{def:g1:plants-around-us:plant}{planta} y su órgano de reserva;
  \item si es materia \omterm{def:g1:animals-around-us:animal}{animal} --- queso, \omterm{def:g2:animals-grow-up:born}{huevo}, jamón --- baja un
        escalón de la cadena: ¿qué comió el \omterm{def:g1:animals-around-us:animal}{animal}? ¿y quién fabricó
        esa comida? sigue hasta la \omterm{def:g1:plants-around-us:plant}{planta};
  \item escribe la cadena con flechas
        (\cref{def:g3:food-chains:chain});
  \item la conclusión es siempre la misma: el rastro termina en una
        \omterm{def:g1:plants-around-us:parts}{hoja}, bajo la \omterm{def:g6:exploring-our-environment:conditions}{luz}.
\end{enumerate}
\end{method}

\section{Ejercicios}

\begin{exercise}[$\star$]\label{exo:g6:plants-are-producers:1}
Define \omterm{def:g6:plants-are-producers:matter}{materia mineral} y \omterm{def:g6:plants-are-producers:matter}{materia orgánica}, con dos ejemplos de cada una.
\end{exercise}

\begin{exercise}[$\star$]\label{exo:g6:plants-are-producers:2}
Enuncia el oficio de las \omterm{def:g5:food-webs:roles}{productoras}: qué entra, de dónde sale la
\omterm{prop:g3:balanced-diet:jobs}{energía}, qué sale y en qué taller ocurre.
\end{exercise}

\begin{exercise}[$\star$]\label{exo:g6:plants-are-producers:3}
Clasifica como mineral u orgánico: el rocío, una \omterm{def:g1:animals-around-us:covering}{concha} de caracol, la
cera de abeja, el granito, el aceite de cocina, la sal marina.
\end{exercise}

\begin{exercise}[$\star$]\label{exo:g6:plants-are-producers:4}
¿De dónde salieron sobre todo los kilos del roble, si apenas salieron
de la tierra?
\end{exercise}

\begin{exercise}[$\star$]\label{exo:g6:plants-are-producers:5}
¿Qué gas toma del aire la \omterm{def:g1:plants-around-us:parts}{hoja} en funcionamiento y cuál suelta? ¿Quién
espira el primero y necesita el segundo?
\end{exercise}

\begin{exercise}[$\star$]\label{exo:g6:plants-are-producers:6}
Nombra cuatro órganos de reserva de las \omterm{def:g1:plants-around-us:plant}{plantas} y qué guarda cada uno,
según el \cref{ex:g6:plants-are-producers:pantry}.
\end{exercise}

\begin{exercise}[$\star\star$]\label{exo:g6:plants-are-producers:7}
¿Por qué se muere de hambre la \omterm{def:g1:plants-around-us:plant}{planta} del armario oscuro con la
despensa de \omterm{prop:g4:what-plants-need:needs}{sales minerales} llena? ¿Qué entrada no tiene sustituto?
\end{exercise}

\begin{exercise}[$\star\star$]\label{exo:g6:plants-are-producers:8}
¿Por qué no puede ser ningún \omterm{def:g1:animals-around-us:animal}{animal} el primer eslabón de una \omterm{def:g3:food-chains:chain}{cadena
alimentaria}? Responde con la distinción entre \omterm{def:g5:food-webs:roles}{productores} y \omterm{def:g5:food-webs:roles}{consumidores}.
\end{exercise}

\begin{exercise}[$\star\star$]\label{exo:g6:plants-are-producers:9}
Pásale el \cref{met:g6:plants-are-producers:trace} a un bocadillo de
queso: pan, mantequilla, queso --- tres rastros hasta tres \omterm{def:g1:plants-around-us:parts}{hojas}.
\end{exercise}

\begin{exercise}[$\star\star$]\label{exo:g6:plants-are-producers:10}
``El compost alimenta a la \omterm{def:g1:plants-around-us:plant}{planta} como el heno alimenta a una vaca''.
Corrige la frase con precisión, usando la
\cref{rem:g4:what-plants-need:soil} y las definiciones de este capítulo
--- ¿qué aporta en realidad el compost y qué fabrica la \omterm{def:g1:plants-around-us:plant}{planta} por sí
misma?
\end{exercise}

\begin{exercise}[$\star\star$]\label{exo:g6:plants-are-producers:11}
Un acuario cerrado e iluminado con \omterm{def:g1:plants-around-us:plant}{plantas} acuáticas y \omterm{prop:g3:sorting-living-things:groups}{peces} puede
funcionar meses; uno cerrado y \emph{a oscuras} fracasa en unos días.
Usa la \cref{rem:g6:plants-are-producers:gas} para explicar los dos
desenlaces.
\end{exercise}

\begin{exercise}[$\star\star\star$]\label{exo:g6:plants-are-producers:12}
Un experimento antiguo: un plantón de sauce criado cinco años en una
tina ganó muchos kilos; la tierra seca de la tina perdió solo unos
gramos. El experimentador concluyó que ``el árbol estaba hecho solo de
agua''. ¿Qué parte de su conclusión era correcta y cuál incompleta ---
y qué segunda entrada mineral, que él no conocía, añadiría hoy una
bióloga?
\end{exercise}

\section{Problema: El año de la huertana}

\begin{problem}[{Problema de fin de semana --- luz, hojas y la
contabilidad de una huerta}]\label{pb:g6:plants-are-producers:1}
Una huertana cultiva patatas, zanahorias, judías y lechugas, y lleva
una contabilidad muy estricta de lo que entra y de lo que sale.

\medskip\noindent\textbf{Parte I --- El acertijo de la contabilidad.}
\begin{enumerate}
  \item La huerta da dos toneladas de verduras al año; la huertana solo
        esparce unos sacos de compost. Pesa el acertijo: ¿de dónde
        salen las toneladas, si no de los sacos? (Responde con el
        oficio de las \omterm{def:g5:food-webs:roles}{productoras}.)
  \item Clasifica la cosecha: la patata, la zanahoria, la \omterm{def:g2:from-seed-to-plant:seed}{semilla} de
        judía, la \omterm{def:g1:plants-around-us:parts}{hoja} de lechuga --- nombra el papel de cada una para
        la \omterm{def:g1:plants-around-us:plant}{planta} (órgano de reserva, órgano en funcionamiento,
        despensa de la \omterm{def:g2:from-seed-to-plant:seed}{semilla}).
  \item La huertana dice: ``mi proveedor de verdad reparte a la \omterm{def:g6:exploring-our-environment:conditions}{luz} del
        día''. Tradúcelo a los términos de la
        \cref{prop:g6:plants-are-producers:produce}.
  \item ¿Para qué necesita entonces el compost la huertana, si no es
        para las toneladas? (Acuérdate de la menta en arena.)
\end{enumerate}

\medskip\noindent\textbf{Parte II --- Leer las estaciones.}
\begin{enumerate}[resume]
  \item Las lechugas de junio, tapadas una semana con una tela de
        prueba, salen pálidas y estiradas. Nombra el experimento que
        repite esto y la entrada que quitó.
  \item Se aporcan las hileras de patatas y se alaban sus matas
        frondosas: ``buenas matas, buenos \omterm{ex:g4:plants-without-seeds:bulbtuber}{tubérculos}''. Explica el
        dicho --- ¿qué relación hay entre las \omterm{def:g1:plants-around-us:parts}{hojas} en funcionamiento
        de arriba y la reserva que engorda abajo?
  \item La sequía de agosto para el crecimiento aunque el sol
        resplandezca. ¿Qué entrada mineral ha fallado y por qué no la
        puede sustituir la \omterm{def:g6:exploring-our-environment:conditions}{luz}?
  \item En septiembre la huertana corta y retira las matas de judías,
        pero entierra sus \omterm{def:g1:plants-around-us:parts}{raíces}. Un vecino lo llama desperdicio.
        Defiende la práctica con las definiciones de materia --- ¿qué
        se le está devolviendo a la tierra y para quién?
\end{enumerate}

\medskip\noindent\textbf{Parte III --- Más allá de la huerta.}
\begin{enumerate}[resume]
  \item El cerdo de la huertana se come las patatas pequeñas; la
        familia se come el cerdo. Escribe la cadena y di dónde se
        fabricó originalmente su \omterm{def:g6:plants-are-producers:matter}{materia orgánica}.
  \item Una clienta afirma que su ensalada es ``sol, sobre todo''.
        Audita la afirmación: enumera los ingredientes reales de la
        \omterm{def:g1:plants-around-us:parts}{hoja} de lechuga y el papel exacto del sol.
  \item Una helada mata de noche las últimas lechugas; en noviembre se
        pudren donde están. Su \omterm{def:g6:plants-are-producers:matter}{materia orgánica} no desaparece.
        Anticipa el \cref{ch:g6:decomposers-and-soil}: ¿quién se la
        queda y hacia qué estado?
  \item Resume la huerta en una sola frase que contenga: \omterm{def:g5:food-webs:roles}{productoras},
        \omterm{def:g6:plants-are-producers:matter}{materia mineral}, \omterm{def:g6:plants-are-producers:matter}{materia orgánica}, \omterm{def:g6:exploring-our-environment:conditions}{luz}.
\end{enumerate}
\end{problem}
